% !PRINTED 777
\weberpage{777}

\begin{center}
II.

\gesperrt{The Irreducibility of the Cyclotomic Equation
and the Theorem on the Primes contained in a
Linear Form.}
\end{center}

In the first supplement, which is to be regarded as a completion of the first Band,
no use was made, of set purpose, of the theory of integral algebraic numbers. In
conclusion we therefore give one more proof of the irreducibility of the cyclotomic
equation, resting on other foundations, namely on the theorems concerning the class
number (\S.\,192); it is remarkable on account of the generalisations which it admits.

Let $m$ be an arbitrary natural number, and let $n$ be the system of the positive
numbers coprime to $m$; among these there are
\weq{1}{\varphi(m) = \mu,}
which are less than $m$, and these we denote by $a$.

If we unite into a single class $A$ all the numbers $n$ that are congruent to one
and the same $a$ with respect to the modulus $m$, we obtain $\mu$ classes of numbers:
\weq{2}{A_1, A_2, \ldots, A_\mu,}
which, under composition by multiplication, form an \gesperrt{Abelian} group of
degree $\mu$, $\mathfrak{N}$. This group has already been considered in \S.\,16 of
this Band.

The $\mu$ characters of this group we denote by
\weq{3}{\chi_1, \chi_2, \ldots, \chi_\mu,}
and we put, when $\chi$ is one of these characters and $n$ denotes a number from
the class $A$,
\weq{4}{\chi(A) = \chi(n).}
These characters, which are more closely determined in \S.\,16, are all
$\mu$th roots of unity.
